\documentclass[10pt,a4paper]{amsart}
\usepackage[margin=19mm]{geometry}
\usepackage{amsmath,amssymb,mathtools}
\usepackage[scr=rsfs,cal=euler]{mathalfa}
\usepackage{xfrac}
\usepackage[unicode,hidelinks,bookmarks=false]{hyperref}
\newcommand{\ie}{i.e.\ }
\newcommand{\cf}{cf.\ }
\newcommand{\resp}{respectively\ }
\newcommand{\Subsection}{\subsection}
\providecommand{\nobreakdash}{\nobreak\hbox{-}}
\setlength{\emergencystretch}{2em}
\multlinegap0pt
\usepackage{mathtools}
\usepackage{amssymb}
\usepackage{enumerate}
\usepackage[shortlabels]{enumitem}
\usepackage{tikz}
\usepackage{float}
\usepackage{xcolor}
\newtheorem{proposition}{Proposition}
\def\Exc{\operatorname{Exc}}
\def\Supp{\operatorname{Supp}}
\title{Coupled Pklt Tuples and Varieties of Pklt Type}
\author{Donghyeon Kim, Dae-Won Lee}
\date{Source: arXiv:2609.11107v1 (CC BY 4.0)}
\begin{document}
\maketitle

\noindent\emph{Source and licence.} Coupled Pklt Tuples and Varieties of Pklt Type. Version arXiv:2609.11107v1 (CC BY 4.0), distributed by arXiv under the Creative Commons Attribution 4.0 International licence, http://creativecommons.org/licenses/by/4.0/, as recorded in the arXiv licence metadata for this exact version and shown on the abstract page https://arxiv.org/abs/2609.11107v1. Full licence notice: \texttt{licenses/arxiv-2609.11107-CC-BY-4.0.txt}.\par\smallskip

\noindent\textit{Editorial scope.} Counted unit: the article's proposition prop:sum (source0.tex lines 411--418). The article's complete original proof of it is delivered with it (source0.tex lines 420--475, one \texttt{proof} environment of 56 lines). No mathematics is written, repaired or re-derived: the statement and the proof are the article's own lines, in the article's own notation, and no retained line is substituted or re-flowed.  No further source-local context is needed: every cross-reference of every retained line resolves inside the two delivered spans.  Nothing was omitted from any retained span, so the package carries no omission record.  No retained line contains a citation, so the wrapper carries no bibliography and no bibliography entry.  No retained line contains a figure environment, an image-inclusion command or drawing code, so no image is delivered and the package declares no asset dependency.  Editorial formatting only, in addition: an AMS \texttt{amsart} preamble that restates the article's own declarations and macros used by the retained lines, namely the environments \texttt{proposition} on separate counters, together with the standard packages those lines need.  The retained environments print in this document as Proposition 1 in order of appearance, whereas the article prints the same items with the numbering of its own full document.  The complete native source of the article as distributed in the arXiv e-print for this exact version is bundled unchanged as \texttt{sources/209988/source0.tex}.
\begin{proposition}\label{prop:sum}
Let $(X,\Delta)$ be a projective klt pair, and let $\mathbf D=(D_1,\ldots,D_r)$ be a tuple of big $\mathbb Q$-Cartier $\mathbb Q$-divisors on $X$. Then
\[
\mathcal J(X,\Delta;\|\mathbf D\|)=\bigcup_{\substack{0\leq D_j'\sim_{\mathbb Q}D_j\\1\leq j\leq r}}
\mathcal J\left(X,\Delta+\sum_{j=1}^rD_j'\right).
\]
More precisely, $\mathcal J(X,\Delta;\|\mathbf D\|)$ is the maximum element.
\end{proposition}

\begin{proof}
Fix effective $\mathbb Q$-divisors $D_j'\sim_{\mathbb Q}D_j$. Choose a sufficiently divisible integer $m>1$ such that every $mD_j'$ is an integral Cartier divisor. Set $H_j\coloneqq mD_j'\in |mD_j|$. Then we have 
\[
\prod_{j=1}^r\mathcal O_X(-H_j)\subseteq \prod_{j=1}^r\mathfrak b(|mD_j|) = \mathfrak a_m(\mathbf D).
\]
The monotonicity of multiplier ideals gives
\begin{align*}
\mathcal J\left(X,\Delta+\sum_{j=1}^rD_j'\right)
&=
\mathcal J\left(X,\Delta;\left(\prod_{j=1}^r\mathcal O_X(-H_j)\right)^{1/m}\right)\\
&\subseteq
\mathcal J\left(X,\Delta;\mathfrak a_m(\mathbf D)^{1/m}\right)\\
&\subseteq
\mathcal J(X,\Delta;\|\mathbf D\|).
\end{align*}

For the reverse inclusion, choose a sufficiently divisible integer
$m>1$ such that
\[
\mathcal J(X,\Delta;\|\mathbf D\|) = \mathcal J\left(X,\Delta;\mathfrak a_m(\mathbf D)^{1/m}\right).
\]
Let \(f\colon Y\to X\) be a common log resolution of $(X,\Delta)$ and the ideals $\mathfrak b(|mD_j|)$. Write \(\mathfrak b(|mD_j|)\mathcal O_Y = \mathcal O_Y(-F_{j,m})\) and \(f^*|mD_j| = |M_{j,m}|+F_{j,m}\), where $|M_{j,m}|$ is basepoint free. Set $F_m\coloneqq\sum_{j=1}^rF_{j,m}$. Then we have
\[
\mathcal J(X,\Delta;\|\mathbf D\|)
= f_*\mathcal O_Y\left(K_Y-\left\lfloor f^*(K_X+\Delta)+\frac{1}{m}F_m\right\rfloor\right).
\]

Choose general members $H_j\in |mD_j|$, and after replacing the resolution if necessary, one may write \(f^*H_j=M_{j,H_j}+F_{j,m}\), where
\[
\Supp\left(f_*^{-1}\Delta+\Exc(f)+F_m+\sum_{j=1}^rM_{j,H_j}\right)
\]
has simple normal crossings. We choose the divisors $M_{j,H_j}$ without common components and without any component in the support of
\[
f^*(K_X+\Delta)+\frac{1}{m}F_m.
\]
Since $m>1$, every component of $M_{j,H_j}$ occurs with coefficient
$\frac{1}{m}<1$. Consequently,
\[
\left\lfloor
f^*(K_X+\Delta)
+\frac{1}{m}F_m
+\frac{1}{m}\sum_{j=1}^rM_{j,H_j}
\right\rfloor
=
\left\lfloor
f^*(K_X+\Delta)+\frac{1}{m}F_m
\right\rfloor.
\]
It follows that
\begin{align*}
\mathcal J(X,\Delta;\|\mathbf D\|)
&=
\mathcal J\left(X,\Delta+\frac{1}{m}\sum_{j=1}^rH_j\right).
\end{align*}
The divisors $\frac{1}{m}H_j$ are effective and $\mathbb Q$-linearly equivalent to $D_j$. The ideal $\mathcal J(X,\Delta;\|\mathbf D\|)$ therefore belongs to the collection.
\end{proof}

\end{document}
